\section{Introduction}
Dual systems express differentiation through rank-one commutators.
For selfadjoint generators, their relation to conjugate variables is part of the analytic approach to free entropy developed by Voiculescu \cite{Voi98}.
Conjugate variables also impose algebraic regularity: they lead to maximal $\Delta$, and maximal $\Delta$ realizes the free field inside the algebra of affiliated operators \cite{ACSY24,MSY23}.
For unbounded dual operators, the domain on which a commutator is known matters when one passes from polynomial differentiation to rational functions.
In particular, an operator initially defined on a polynomial algebra need not have all of $\M\Omega$ in its domain.

This paper compares the actual adjoint condition for a derivation on a represented star algebra with dual identities on $\M\Omega$.
For bounded normal generators the two conditions are equivalent when the dual operators and their Hilbert space adjoints are both required to be defined on $\M\Omega$.
The main analytic step is a uniform estimate for both slices of a coordinate derivation, without a freeness assumption on the generators.
It applies to the canonical polynomial operators themselves and therefore identifies the closures used in the resulting rationality criterion.
The comparison changes for nonnormal generators: we construct examples separating the two hypotheses in both directions.

Throughout, $(\M,\tau)$ is a finite von Neumann algebra with a faithful normal probability trace,
\[
 H:=L^2(\M,\tau),\qquad \Omega:=1,\qquad
 \langle a\Omega,b\Omega\rangle:=\tau(a^*b).
\]
Thus the inner product is linear in its second variable.
We identify $a\in\M$ with $a\Omega$ when no domain ambiguity results.
Let $X=(X_1,\ldots,X_d)$ be a finite tuple of bounded operators generating $\M$ as a von Neumann algebra, and set
\[
 A_X:=\C\langle X_1,\ldots,X_d\rangle,\qquad
 \A_X:=\C\langle X_1,\ldots,X_d,X_1^*,\ldots,X_d^*\rangle.
\]
These are represented algebras; their relations are retained.
In particular, the adjoints are not independent formal generators.
The outer bimodule actions on the Hilbert space tensor product $H\bar\otimes H$ are
\[
 a\cdot(\eta\otimes\theta)\cdot b:=(a\eta)\otimes(\theta b).
\]

\begin{definition}\label{def:hypotheses}
We use the following three conditions on $X$.
The condition $\mathsf A(X)$ means that there are complex-linear outer derivations
\[
 \partial_i:\A_X\Omega\subset H\longrightarrow H\bar\otimes H,
 \qquad \partial_iX_j=\delta_{ij}\Omega\otimes\Omega,
 \qquad \Omega\otimes\Omega\in\Dom\partial_i^*.
\]
Their conjugate vectors are $\xi_i:=\partial_i^*(\Omega\otimes\Omega)$.
The condition $\mathsf S(X)$ means that there are densely defined complex-linear operators $D_i$ on $H$ such that
\[
 \M\Omega\subseteq\Dom D_i\cap\Dom D_i^*,\qquad
 [X_j,D_i]a\Omega=\delta_{ij}P_\Omega a\Omega\quad(a\in\M).
\]
The condition $\mathsf F(X)$ requires only complex-linear maps $D_i:\M\Omega\to H$ with these commutator identities.
It imposes no continuity, closability or adjoint-domain requirement.
\end{definition}
Here and below $P_\Omega$ denotes the orthogonal projection onto $\C\Omega$, so $P_\Omega a\Omega=\tau(a)\Omega$.
Every condition in Definition~\ref{def:hypotheses} is existential.
The condition $\mathsf A$ is the precise starred-core hypothesis in \cite[Proposition~3.5]{ACSY24}; no reality condition on $\partial_i$ is part of it.
Plainly $\mathsf S$ implies $\mathsf F$.

\begin{theorem}\label{thm:normal}
Suppose every $X_j$ is a bounded normal operator.
Then $\mathsf A(X)$ and $\mathsf S(X)$ are equivalent.
More precisely, under $\mathsf A(X)$ there is an absolute constant $C$ such that
\begin{equation}\label{eq:main-slices}
 \| (\id\otimes\tau)\partial_i p\|_2+
 \| (\tau\otimes\id)\partial_i p\|_2
 \le C\|\xi_i\|_2\|p\|\qquad(p\in\A_X).
\end{equation}
Define $E_i^\circ(p\Omega):=(\id\otimes\tau)\partial_i p$ on exactly $\A_X\Omega$ and $D_i:=-\overline{E_i^\circ}$.
Then
\[
 \M\Omega\subseteq\Dom D_i\cap\Dom D_i^*,\qquad
 [X_j,D_i]|_{\M\Omega}=\delta_{ij}P_\Omega|_{\M\Omega}.
\]
Conversely, every system witnessing $\mathsf S(X)$ determines the same starred derivations, with
\begin{equation}\label{eq:strong-vacuum}
 \xi_i=-D_i^*\Omega+(D_i\Omega)^*.
\end{equation}
\end{theorem}

The sign in the canonical choice is dictated by the generator-first convention for commutators.
Theorem~\ref{thm:normal} extends the polynomial-core construction of \cite[Theorem~3.4]{ACSY24} by giving both actual domains on $\M\Omega$.
Its proof uses the normal spectral phase, partial liberation, and a signed curvature estimate.
An estimate first proved for separated spectral components is transferred to arbitrary spectra through finite lattice approximations and a projection encoding of the other coefficients.

Write $L^0(\M)$ for the algebra of affiliated operators, with sums and products understood as affiliated closures.
An element of $\M$ is called regular if it is invertible in $L^0(\M)$, equivalently if its action on $H$ is injective with dense range.

\begin{definition}\label{def:rational}
Put $\mathcal D_X:=\Div_{L^0(\M)}A_X$, the smallest unital complex subalgebra containing $A_X$ and closed under inverses of its elements that are invertible in $L^0(\M)$.
For $u\in L^0(\M)$, a bounded regular fraction is a quadruple $Q=(s,a,b,t)\in\M^4$ with $s,t$ regular and
\[
 u=s^{-1}a=bt^{-1}.
\]
For a specified operator $D$, define
\[
 \Gamma[Q,D]:=sDb-aDt,
 \qquad \Dom\Gamma[Q,D]:=\{v\in H:bv,tv\in\Dom D\}.
\]
Finite rank of this expression means finite-dimensional range on its entire stated domain.
\end{definition}

Let $\rho(x)v:=vx$ denote bounded right multiplication when $x\in\M$, and let $\F(H)$ be the finite-rank bounded operators on $H$.
\begin{theorem}\label{thm:rational}
Suppose $X$ is a normal generating tuple satisfying either equivalent condition of Theorem~\ref{thm:normal}, and let $D_i$ be its canonical generator-first operators.
For every $u\in L^0(\M)$ the following are equivalent:
\begin{enumerate}
\item $u\in\mathcal D_X$;
\item there is one bounded regular fraction $Q=(s,a,b,t)$ for $u$ and a tuple $T_i\in\F(H)$ satisfying
\begin{equation}\label{eq:boundary}
 \sum_i[\rho(X_i),T_i]=aP_\Omega t-sP_\Omega b;
\end{equation}
\item for one common bounded regular fraction $Q$ for $u$, every $\Gamma[Q,D_i]$ has finite rank;
\item for every bounded regular fraction $Q$ for $u$, every $\Gamma[Q,D_i]$ has finite rank.
\end{enumerate}
In the finite-rank case these expressions have bounded finite-rank closures, and
\[
 T_i=-\overline{\Gamma[Q,D_i]}
\]
is the unique solution of \eqref{eq:boundary} at that fraction.
The numerical rank in each coordinate is independent of $Q$.
For any supplied system witnessing $\mathsf S(X)$, its fraction expressions agree with the canonical ones on $\M\Omega$ and have the same bounded finite-rank closures whenever either has finite rank on its natural domain.
\end{theorem}

The algebra in this theorem is generated by the unstarred coordinates.
The star algebra supplies the analytic core, rather than changing the meaning of rationality.
The existential fraction condition in item~2 also holds for arbitrary bounded tuples with maximal $\Delta$; its specified-operator realization uses the additional normal and domain hypotheses.

\begin{theorem}\label{thm:general}
For arbitrary finite bounded generating tuples, $\mathsf A$ and $\mathsf F$ are incomparable, and $\mathsf A$ and $\mathsf S$ are incomparable.
There is a two-coordinate tuple satisfying $\mathsf A$ and failing $\mathsf F$, and a one-coordinate tuple satisfying $\mathsf S$ for which no normalized Hilbert-valued derivation on the represented star algebra exists.
Either $\mathsf A(X)$ or $\mathsf F(X)$ implies $\Delta(X)=d$.
\end{theorem}

We use $\Delta$ in the sense of \cite{MSY23} and only need its maximality characterization:
\begin{equation}\label{eq:cycles}
 \Delta(X)=d\quad\Longleftrightarrow\quad
 \left[\sum_i[\rho(X_i),T_i]=0,\quad T_i\in\F(H)\right]
 \Longrightarrow T_1=\cdots=T_d=0.
\end{equation}
This characterization and the free-field realization are valid without normality \cite[Theorems~3.1 and~1.1]{MSY23}.
The implication from $\mathsf A$ is \cite[Proposition~3.5]{ACSY24}.
We give a direct proof for the algebraic hypothesis $\mathsf F$ below.
The final circular example shows that even $\mathsf A$ together with a bounded canonical dual does not extend Theorem~\ref{thm:rational} to arbitrary nonnormal generators.
